\section{Beam 2: competence evolution through experience}

Beam~2 is not taken from a single reference. Its components are supported by complementary
literatures.

\subsection{Learning by doing}

The broad principle that productive activity can itself generate learning is classical.
\citet{arrow1962} formalized learning by doing as an economic mechanism through which
experience accumulated in production affects future productivity. For the present purpose,
the important proposition is modest: operation need not only consume current competence;
it can also contribute to future capability.

\subsection{Task specificity}

A generic stock of experience is too coarse for the intended interface with heterogeneous
problems. The task-specific human-capital literature, including \citet{gibbonswaldman2004},
supports the idea that experience accumulated in particular tasks can produce capability
that is especially relevant to those tasks.

This motivates indexing experience and competence by a task or problem region \(q\in\Q\),
without yet prescribing a particular representation of \(\Q\).

\subsection{Organizational learning: experience to knowledge}

\citet{argotemiron2011} provide the most direct conceptual bridge. They define organizational
learning around a change in organizational knowledge that occurs as a function of experience.
Their framework represents an ongoing cycle in which task-performance experience is converted
into knowledge; knowledge becomes embedded in organizational context; and the changed context
affects future experience.

The framework distinguishes three learning subprocesses:
\begin{enumerate}
\item knowledge creation from experience;
\item knowledge retention; and
\item knowledge transfer.
\end{enumerate}

For the present two-beam skeleton, creation and retention are already sufficient. Explicit
transfer can be added later and is not required to make the current feedback loop coherent.

Argote and Miron-Spektor also emphasize that experience should be characterized at a
fine-grained level. This is compatible with treating the experience produced by an allocation
as task-dependent rather than as a single undifferentiated count.

Their framework is conceptual rather than an optimization model. It does not provide a unique
state-transition equation for HLS.

\subsection{An explicit dynamic competence model}

As a primary current formal reference for this beam, \citet{gutjahr2011} provides a
quantitative model in which current work allocation changes future competence.

For project/task class \(i\) and period \(t\), let \(z_{it}\) be a competence score and
\(x_{it}\) the share of available work capacity invested in class \(i\). Gutjahr models
competence as
\begin{equation}
z_{it}
=
z_{i1}
-
\beta_i(t-1)
+
\eta_i\sum_{s=1}^{t-1}x_{is},
\label{eq:gutjahr}
\end{equation}
where \(\eta_i\) represents learning and \(\beta_i\) knowledge depreciation.

Competence is mapped to productive efficiency through a nondecreasing learning function
\begin{equation}
\gamma_{it}=\phi(z_{it}),\qquad 0\leq\gamma_{it}\leq 1.
\end{equation}
Gutjahr gives a concrete operational meaning to \(\gamma_{it}\): if a perfectly competent
team requires \(d\) units of real work time, a team with efficiency \(\gamma_{it}\) requires
\(d/\gamma_{it}\).

The essential dynamic mechanism is therefore
\begin{equation}
\text{work allocation}
\longrightarrow
\text{experience}
\longrightarrow
\text{competence}
\longrightarrow
\text{future productive efficiency}.
\end{equation}

The paper also demonstrates that the optimal pattern of concentration or mixing depends on
the learning function, depreciation and horizon. Thus learning does not mechanically imply
either specialization or diversification.

\subsection{What Beam 2 establishes for the present skeleton}

The combined literature supports the following ingredients:
\begin{enumerate}
\item task performance can generate learning;
\item learning can be task-specific;
\item knowledge/competence can be retained and can depreciate;
\item current work allocation can therefore alter future productive capability.
\end{enumerate}

A generic HLS-oriented transition can consequently be written as
\begin{equation}
C_{t+1}=\mathcal{L}(C_t,e_t;\theta_t),
\label{eq:learning}
\end{equation}
where \(e_t\) is learning-relevant exposure and \(\theta_t\) contains learning and retention
parameters. Real work and exposure coincide in Gutjahr's source model through its stated
work-dependent transition; an HLS translation must preserve or explicitly replace that
mapping. Equation~\eqref{eq:learning} is a synthesis, not a formula asserted by the cited
literature.

\subsection{What Beam 2 does not establish}

Gutjahr's competence is aggregated at team level in the basic model. It is not a theory of
routing among heterogeneous machine-learning models. Nor does the cited foundation determine
how knowledge should be transferred from a teacher model to a learner, how machine-learning
competence should be measured, or how experience generalizes across neighboring task regions.

Those questions lie outside the present note.

Gutjahr is not the theory of Beam~2. Its formalization is one source among multiple
independent lines, and its team-level state, work variable, and efficiency mapping must not
be imposed on HLS without an explicit ontology mapping.
